\section{Clôture algébrique du mode de Fourier nul associé à la périodicité CT108 dans la fibre électronique}
\label{sec:clausura-nulo-material-electron-rev3}
\index{électron!clôture nulle--matérielle}
\index{CT108!mode nul et électron}

L'identité autoduale de Planck fixe l'origine de la droite positive des
masses ; la généalogie électronique fixe son premier déplacement matériel
complet. Cette section réunit, sans répéter leurs démonstrations de référence,
la réalisation circulaire de CT108, la composition basale de l'électron et
l'holonomie bidirectionnelle de sa voie.

Soit
\[
 \mathcal H_{108}^{\rm reg}
 :=\ell^2(\mathbb Z/108\mathbb Z;\mathbb R),
 \qquad
 S_{108}^{\rm reg}\delta_\theta=\delta_{\theta+1},
\]
la représentation régulière réelle du décalage cyclique de CT108, et soit
\[
 n_0:=\frac1{\sqrt{108}}\sum_{\theta=0}^{107}\delta_\theta,
 \qquad
 \mathcal N_{108}:=\ker(S_{108}^{\rm reg}-I)=\mathbb R n_0.
\]
Le mode de Fourier nul normalisé est ainsi défini, et non seulement nommé.
Soit \(\mathcal G_e\) l'espace des généalogies électroniques enrichies de
\cref{prop:planck-ct108-realizacion-centro-electron-rev3} et fixons le
sous-module linéaire
\(\ell_{\Gamma_e}:=\mathbb R g_{\Gamma_e}
\subset\mathbb R[\mathcal G_e]\) engendré par la voie centrale dans la
linéarisation libre de cet espace. De même, si
\(F_e^+=\mathbb R_{>0}I\) est la demi-droite positive déjà définie, soit
\(\mathcal F_e:=\operatorname{span}_{\mathbb R}(F_e^+)=\mathbb R I\) son
enveloppe linéaire, et soit
\(\mathcal L_M\) la droite de masse définie dans
\cref{sec:ley-general-masa-accion-autonoma}. Les bases déclarées déterminent
le morphisme linéaire de rang un
\begin{equation}
 \mathfrak C_e:
 \mathcal N_{108}\otimes\ell_{\Gamma_e}
 \longrightarrow \mathcal F_e\otimes\mathcal L_M,
 \qquad
 \mathfrak C_e(n_0\otimes g_{\Gamma_e})
 =\mathfrak c_e\,I\otimes\mathcal U_M,
 \label{eq:morfismo-clausura-nulo-material-electron-rev3}
\end{equation}
où la coordonnée positive est
\begin{equation}
 \boxed{
 \mathfrak c_e=
 \frac{\sqrt3}{4}
 \left[
  \exp\!\left(\frac{\varphi_{\rm HMT}}{\pi_{\rm HMT}^{2}}\right)
  -22\alpha_{\rm HMT}^{3}
 \right]
 (1+15\Delta_4)
 R_{\rm act}^{\,80-54\alpha_{\rm HMT}+6\Delta_4}.}
 \label{eq:coordenada-clausura-nulo-material-electron-rev3}
\end{equation}

\begin{theorem}[Clôture électronique nulle--matérielle de rang un]
\label{thm:clausura-nulo-material-electron-rev3}
Avec les primitives et les cartes déjà fixées, le morphisme
\(\mathfrak C_e\) factorise exactement la sortie électronique bêta. Ses
quatre facteurs conservent, dans cet ordre, l'incompatibilité somme--produit
d'APP ; l'orientation aréale--volumétrique avec le déficit de frontière ;
la torsion et la mémoire retenue ; et l'holonomie bidirectionnelle de la voie
électronique. Dans la carte \(\mathcal U_M=1\,\mathrm{MeV}/c^2\),
\begin{equation}
 \mathfrak C_e(n_0\otimes g_{\Gamma_e})
 =0.5109989506900008086082571648\ldots\,
  I\otimes\mathrm{MeV}/c^2.
 \label{eq:masa-clausura-nulo-material-electron-rev3}
\end{equation}
\end{theorem}

\begin{proof}
Le \cref{thm:cierre-electron-holografico-rev6} démontre que les trois
premiers facteurs de
\eqref{eq:coordenada-clausura-nulo-material-electron-rev3} composent le
scalaire basal positif \(\mathfrak e_{\rm HMT}\). Le
\cref{thm:reduccion-electron-two-way-rev11} obtient, à partir de deux lecteurs
indépendants de la voie,
\[
 K_D(\Gamma_e)=K_\Omega(\Gamma_e)=(80,54,6),
 \qquad
 \kappa_e=80-54\alpha_{\rm HMT}+6\Delta_4.
\]
Le rapport \(R_{\rm act}>0\) appartient à l'unique droite d'action
pré-retour/post-retour. Par conséquent, le quatrième facteur est positif et
la composition coïncide avec
\eqref{eq:electron-beta-two-way-rev9} et
\eqref{eq:masa-electron-two-way-rev11}. La carte dimensionnelle n'agit
qu'après cette factorisation.
\end{proof}

\paragraph{Portée mathématique et réalisation physique.}
Le résultat exact est la réalisation linéaire déclarée de la factorisation
électronique sur le mode nul, la généalogie et la fibre typés. Le mode
nul de CT108 ne sélectionne pas à lui seul la coordonnée \(\mathfrak c_e\) :
celle-ci provient des quatre facteurs démontrés à leurs emplacements et le
morphisme les réunit en une unique clôture algébrique de rang un. La dynamique électromagnétique
liée conserve comme exigences propres un hamiltonien stable, la charge,
le spin \(1/2\), le facteur de forme et le couplage de jauge ; ces objets ne
sont pas des prémisses de la factorisation précédente. L'unité de Planck occupe le
point autodual de la droite et l'électron occupe une section généalogiquement
décalée : une interprétation comme trou noir classique appartient à un autre
type de réalisation et n'intervient pas dans ce morphisme.
